\begin{statementbox}
\begin{center}\begin{tikzpicture}[font=\small,>=stealth,line join=round]
\draw[gray!40,rounded corners=3pt] (0,0) rectangle (14.2,8.3);\node at (7.1,7.700000000000001) {一次互为倒数，两次回到同一高度};
\begin{scope}[shift={(1.7923076923076926,2.0951612903225807)},x=1.153846153846154cm,y=1.4838709677419355cm]\draw[->,gray] (-0.6,0)--(4.6,0) node[right]{$x$};\draw[->,gray] (0,-0.3)--(0,2.8) node[above]{$y$};\draw[blue,thick,solid,domain=-0.3:4.3,samples=180] plot(\x,{2^(sin(deg(pi*\x)))});\fill[orange!85!black] (0.5,2) circle(2.2pt);\fill[orange!85!black] (1.5,0.5) circle(2.2pt);\fill[orange!85!black] (2.5,2) circle(2.2pt);\draw[gray,dashed] (-.5,1)--(4.4,1);\node[left] at (0,1){$1$};\foreach \v/\lab in {.5/\frac12,1.5/\frac32,2.5/\frac52}{\node[below=5pt] at (\v,0){$\lab$};}\end{scope}\node[anchor=west,align=left,text width=5.0cm] at (8.4,4.0) {$f=2^{\sin(\pi x)},\ a=1$\\[9pt]$f(\frac12)=2$\\[9pt]$f(\frac32)=\frac12$\\[9pt]$f(\frac52)=2$\\[9pt]正周期可取2};
\node[text width=12.6cm,align=center] at (7.1,.65) {三点高度是2、二分之一、2；倒数不同于把高度变为相反数。};
\end{tikzpicture}\end{center}


\textbf{\textcolor{CStatement}{先说它能解决什么问题}}

题目没有告诉你$f(x)$究竟是什么函数，只告诉你：\textbf{向右走一段，函数值就变成原来的倒数（或负倒数）}。这时不必先猜函数表达式，只要继续走同样长的一段，就能发现函数值又回来了，从而迅速判断周期，并计算很远处的函数值。

\medskip
\textbf{\textcolor{CStatement}{结论一：变倒数，再变一次就回来了}}

设$f(x)$是定义在\textbf{全体实数}上的实值函数，且$f(x)\neq0$。设$a\neq0$为固定常数。如果下列关系中的\textbf{某一种在所有$x$处始终成立}：
\[
\boxed{\ f(x+a)=\frac{1}{f(x)}\ }
\qquad\text{或}\qquad
\boxed{\ f(x+a)=-\frac{1}{f(x)}\ },
\]
那么
\[
\boxed{\ f(x+2a)=f(x)\ },\qquad
\boxed{\ T=2|a|\ \text{是一个正周期}}.
\]
如果题目已经给出$a>0$，可以直接写$T=2a$。\textbf{这里说的是“一个周期”，不是必然的“最小正周期”。}

\medskip
\end{statementbox}\clearpage\begin{statementbox}\textbf{\textcolor{CStatement}{结论二：比“倒数”更一般}}

只要$k$是固定的\textbf{非零实数}，且
\[
\boxed{\ f(x+a)=\frac{k}{f(x)}\ },\quad f(x)\neq0,\quad a\neq0,
\]
同样有
\[
\boxed{\ f(x+2a)=f(x)\ },\qquad \boxed{\ T=2|a|\text{是一个周期}}.
\]
因为连续两次变换是
\[
\underbrace{y\ \longrightarrow\ \frac{k}{y}\ \longrightarrow\
\frac{k}{k/y}=y}_{\text{两次后回到原数}}\qquad (k\neq0,\ y\neq0).
\]
前面两种情况，只是这里的$k=1$和$k=-1$。

\medskip
\textbf{\textcolor{CStatement}{结论三：以后还会重复}}
\[
 f(x+2na)=f(x)\quad(n\in\mathbb Z).
\]
这句话表示：每隔$2|a|$就回到同样的函数值，向左、向右都一样。

\medskip
\textbf{\textcolor{CStatement}{条件提醒}}

如果函数只在某个范围内有定义，就要先确认$x$、$x+a$、$x+2a$等需要用到的位置都在定义域内。要把$2|a|$称为整个函数的周期，还应保证定义域整体向左或向右移动$2|a|$后不变。例如，定义域为$\mathbb R$就没有这个顾虑。
\end{statementbox}
